% ============================================================
\chapter{Apprentissage Machine en Finance}
\label{ch:ml_finance}
% ============================================================

\section{Introduction}

L'apprentissage machine\index{machine learning!finance} (ML) transforme la
finance quantitative en permettant de d\'ecouvrir des structures non lin\'eaires
dans les donn\'ees de march\'e. Ce chapitre pr\'esente les principales
applications : trading algorithmique, pricing d'options par r\'eseaux de
neurones, gestion de portefeuille par apprentissage par renforcement, analyse
de sentiment par NLP, ainsi que les risques et la r\'eglementation.

\begin{intuition}
Les mod\`eles classiques (Black-Scholes, GARCH) imposent une structure
param\'etrique. Le ML permet de \og laisser parler les donn\'ees \fg en
apprenant des relations complexes directement, au prix d'un risque de
surapprentissage accru et d'une moindre interpr\'etabilit\'e.
\end{intuition}

\section{Trading algorithmique}

\begin{definition}[Strat\'egie de trading algorithmique]
\index{trading algorithmique}
Une \textbf{strat\'egie de trading algorithmique} est une r\`egle syst\'ematique
$\delta_t = g(x_t; \theta)$ qui d\'etermine la position $\delta_t$ \`a
l'instant $t$ en fonction de features observables $x_t$ et de param\`etres
$\theta$ optimis\'es sur donn\'ees historiques.
\end{definition}

\begin{proposition}[Features classiques]
Les pr\'edicteurs typiques incluent :
\begin{enumerate}
    \item \textbf{Indicateurs techniques} : moyennes mobiles, RSI, bandes de
    Bollinger, MACD.
    \item \textbf{Microstructure} : d\'es\'equilibre du carnet d'ordres (order
    flow imbalance), spread bid-ask, volume.
    \item \textbf{Features fondamentaux} : ratios financiers, earnings
    surprises.
    \item \textbf{Features alternatifs} : donn\'ees satellites, trafic web,
    sentiment de r\'eseaux sociaux.
\end{enumerate}
\end{proposition}

\begin{attention}[title=\textbf{Surapprentissage et data snooping}]
Le backtesting massif sur donn\'ees historiques produit des strat\'egies qui
performent parfaitement in-sample mais \'echouent out-of-sample. Il faut :
validation crois\'ee temporelle (walk-forward), correction de Bonferroni pour
les tests multiples, et r\'eserve de donn\'ees fraïches pour le test final.
\end{attention}

\begin{theoreme}[Borne de g\'en\'eralisation]
\index{borne de g\'en\'eralisation}
Pour une classe d'hypoth\`eses $\mathcal{H}$ de dimension de Vapnik-Chervonenkis
$d_{\mathrm{VC}}$, avec probabilit\'e au moins $1 - \delta$ sur un \'echantillon
de taille $n$ :
\[
    R(h) \leq \hat{R}(h) + \sqrt{\frac{d_{\mathrm{VC}}(\ln(2n/d_{\mathrm{VC}}) + 1) + \ln(4/\delta)}{n}},
\]
o\`u $R(h)$ est le risque r\'eel et $\hat{R}(h)$ le risque empirique.
\end{theoreme}

\section{R\'eseaux de neurones pour le pricing d'options}

\begin{definition}[R\'eseau de neurones feedforward]
\index{r\'eseau de neurones}
Un \textbf{r\'eseau de neurones feedforward} \`a $L$ couches est une fonction
\[
    f_\theta(x) = W_L\,\sigma(W_{L-1}\,\sigma(\cdots\sigma(W_1 x + b_1)\cdots) + b_{L-1}) + b_L,
\]
o\`u $\sigma$ est une fonction d'activation non lin\'eaire (ReLU, tanh, etc.),
$W_\ell$ les matrices de poids et $b_\ell$ les biais.
\end{definition}

\begin{theoreme}[Approximation universelle]
\index{approximation universelle}
Pour toute fonction continue $g : \R^d \to \R$ sur un compact $K$ et tout
$\epsilon > 0$, il existe un r\'eseau feedforward \`a une couche cach\'ee avec
activation sigmo\"ide tel que $\sup_{x \in K}\abs{f_\theta(x) - g(x)} < \epsilon$.
\end{theoreme}

\begin{proposition}[Pricing par r\'eseau de neurones]
L'approche consiste \`a :
\begin{enumerate}
    \item G\'en\'erer $N$ jeux de param\`etres $(S_0^{(i)}, K^{(i)}, T^{(i)},
    \sigma^{(i)}, r^{(i)})$ et calculer les prix $V^{(i)}$ par Black-Scholes
    ou Monte Carlo.
    \item Entra\^iner un r\'eseau $f_\theta$ minimisant
    $\frac{1}{N}\sum_{i=1}^N (f_\theta(x^{(i)}) - V^{(i)})^2$.
    \item En inf\'erence, $f_\theta$ donne le prix et les grecs
    ($\Delta = \partial f_\theta/\partial S$) quasi instantan\'ement.
\end{enumerate}
\end{proposition}

\begin{remarque}
Les r\'eseaux de neurones acc\'el\`erent le pricing d'un facteur $10^3$ \`a
$10^5$ par rapport \`a Monte Carlo, ce qui est crucial pour le calcul de la
XVA (CVA, DVA, FVA) sur des portefeuilles contenant des milliers de
transactions.
\end{remarque}

\section{Apprentissage par renforcement pour la gestion de portefeuille}

\begin{definition}[Processus de d\'ecision markovien (MDP)]
\index{MDP}
Un \textbf{MDP} est un tuple $(\mathcal{S}, \mathcal{A}, P, R, \gamma)$ o\`u
$\mathcal{S}$ est l'espace d'\'etats, $\mathcal{A}$ l'espace d'actions,
$P(s'|s, a)$ la dynamique de transition, $R(s, a)$ la r\'ecompense, et
$\gamma \in [0, 1)$ le facteur d'actualisation.
\end{definition}

\begin{definition}[Politique et fonction de valeur]
\index{politique!apprentissage par renforcement}
Une \textbf{politique} $\pi : \mathcal{S} \to \mathcal{A}$ sp\'ecifie l'action
\`a chaque \'etat. La \textbf{fonction de valeur} est
$V^\pi(s) = \E_\pi\left[\sum_{t=0}^\infty \gamma^t R(s_t, a_t) \mid s_0 = s\right]$.
La politique optimale satisfait $V^*(s) = \max_\pi V^\pi(s)$.
\end{definition}

\begin{theoreme}[\'Equation de Bellman]
\index{\'equation de Bellman}
La fonction de valeur optimale satisfait
\[
    V^*(s) = \max_{a \in \mathcal{A}} \left[R(s, a)
    + \gamma \sum_{s'} P(s'|s, a)\,V^*(s')\right].
\]
\end{theoreme}

\begin{proposition}[Application \`a la gestion de portefeuille]
\begin{itemize}
    \item \textbf{\'Etat} : $s_t = (w_t, S_t, \text{features}_t)$ (richesse,
    prix, indicateurs).
    \item \textbf{Action} : $a_t = (\delta_1, \ldots, \delta_d)$ (allocations
    dans $d$ actifs).
    \item \textbf{R\'ecompense} : $R_t = \ln(w_{t+1}/w_t)$ (rendement
    logarithmique) ou ratio de Sharpe diff\'erentiel.
    \item \textbf{Algorithmes} : DDPG, PPO, SAC pour actions continues.
\end{itemize}
\end{proposition}

\section{NLP et analyse de sentiment}

\begin{definition}[Analyse de sentiment financi\`ere]
\index{analyse de sentiment}
L'\textbf{analyse de sentiment} en finance consiste \`a extraire un signal
haussier/baissier \`a partir de sources textuelles : communiqu\'es de presse,
rapports d'analystes, tweets, transcriptions d'appels de r\'esultats.
\end{definition}

\begin{proposition}[Pipeline NLP pour la finance]
\begin{enumerate}
    \item \textbf{Collecte} : flux Reuters/Bloomberg, API Twitter/Reddit.
    \item \textbf{Pr\'etraitement} : tokenisation, suppression des stop words,
    lemmatisation.
    \item \textbf{Mod\`ele} : FinBERT (BERT fine-tun\'e sur du texte financier)
    ou GPT sp\'ecialis\'e.
    \item \textbf{Agr\'egation} : score de sentiment moyen pond\'er\'e par
    source et par r\'ecence.
    \item \textbf{Signal de trading} : combinaison avec des features quantitatifs.
\end{enumerate}
\end{proposition}

\begin{formules}[title=\textbf{M\'etriques de performance des strat\'egies ML}]
\begin{align*}
    \text{Ratio de Sharpe} &= \frac{\E[R_p] - r_f}{\sqrt{\Var(R_p)}} \\
    \text{Ratio de Sortino} &= \frac{\E[R_p] - r_f}{\sqrt{\E[\min(R_p - r_f, 0)^2]}} \\
    \text{Maximum Drawdown} &= \max_{t} \frac{\max_{s \leq t} W_s - W_t}{\max_{s \leq t} W_s}
\end{align*}
\end{formules}

\section{Risques et r\'eglementation}

\begin{definition}[Risque de mod\`ele]
\index{risque de mod\`ele}
Le \textbf{risque de mod\`ele} en ML inclut :
\begin{itemize}
    \item \textbf{Surapprentissage} : le mod\`ele m\'emorise le bruit.
    \item \textbf{Distribution shift} : les donn\'ees futures diff\`erent des
    donn\'ees d'entra\^inement (changements de r\'egime).
    \item \textbf{Bo\^ite noire} : impossibilit\'e d'expliquer les d\'ecisions
    (probl\`eme d'interpr\'etabilit\'e).
    \item \textbf{Attaques adversariales} : manipulation d\'elib\'er\'ee des
    inputs pour tromper le mod\`ele.
\end{itemize}
\end{definition}

\begin{remarque}
Le r\`eglement europ\'een AI Act (2024) classe les syst\`emes de scoring de
cr\'edit et de trading algorithmique comme \`a \og haut risque \fg, imposant
des exigences de transparence, de documentation et de supervision humaine.
\end{remarque}

\begin{codeblock}[title=\textbf{Pricing d'options par r\'eseau de neurones}]
\begin{minted}{python}
import numpy as np
import torch
import torch.nn as nn
from scipy.stats import norm

def bs_call(S, K, T, r, sigma):
    d1 = (np.log(S/K) + (r + sigma**2/2)*T) / (sigma*np.sqrt(T))
    d2 = d1 - sigma*np.sqrt(T)
    return S*norm.cdf(d1) - K*np.exp(-r*T)*norm.cdf(d2)

# Generate training data
n_train = 50_000
S = np.random.uniform(50, 150, n_train)
K = np.random.uniform(50, 150, n_train)
T = np.random.uniform(0.1, 2.0, n_train)
sigma = np.random.uniform(0.1, 0.5, n_train)
r_rate = np.full(n_train, 0.05)
prices = bs_call(S, K, T, r_rate, sigma)

X = torch.tensor(np.column_stack([S, K, T, sigma]), dtype=torch.float32)
y = torch.tensor(prices, dtype=torch.float32).unsqueeze(1)

model = nn.Sequential(
    nn.Linear(4, 64), nn.ReLU(),
    nn.Linear(64, 64), nn.ReLU(),
    nn.Linear(64, 1))
optimizer = torch.optim.Adam(model.parameters(), lr=1e-3)

for epoch in range(200):
    pred = model(X)
    loss = nn.MSELoss()(pred, y)
    optimizer.zero_grad(); loss.backward(); optimizer.step()
print(f"MSE finale: {loss.item():.6f}")
\end{minted}
\end{codeblock}

\section{Exercices}

\begin{exercice}[Feature engineering]
Pour une strat\'egie de trading sur le S\&P 500, proposer 10 features
pertinents. Discuter les probl\`emes de look-ahead bias et de
multicollin\'earit\'e.
\end{exercice}

\begin{exercice}[R\'eseau pour les grecs]
Entra\^iner un r\'eseau de neurones pour pricer des calls europ\'eens. Calculer
le delta par diff\'erentiation automatique et comparer avec la formule
analytique.
\end{exercice}

\begin{exercice}[RL pour allocation]
Impl\'ementer un agent DQN simple pour l'allocation entre un actif risqu\'e
(rendement $\mathcal{N}(8\%, 20\%)$) et un actif sans risque (3\%). Comparer
la politique apprise avec la r\`egle de Merton.
\end{exercice}

\begin{exercice}[Analyse de sentiment]
Collecter 1000 titres d'actualit\'es financi\`eres et les classifier en
positif/n\'egatif/neutre. Mesurer la corr\'elation entre le score de sentiment
agr\'eg\'e et les rendements du lendemain.
\end{exercice}

\begin{exercice}[Validation temporelle]
Expliquer pourquoi la validation crois\'ee standard (k-fold) est
inappropri\'ee en finance. Impl\'ementer une validation walk-forward avec une
fen\^etre glissante de 2 ans d'entra\^inement et 6 mois de test.
\end{exercice}

\begin{exercice}[Interpr\'etabilit\'e]
Appliquer les valeurs de Shapley (SHAP) \`a un mod\`ele XGBoost de pr\'ediction
de rendements. Identifier les features les plus importants et discuter la
stabilit\'e des explications au cours du temps.
\end{exercice}
